\section{An ensemble benchmark derived by marginalization}\label{app:ensemble}

To distinguish a model joint law from the canonical construction,
consider the circular density
\begin{equation}\label{eq:cue-joint}
 p_N^{\rm circ}(\boldsymbol\theta)
   =\frac1{N!(2\pi)^N}
      \prod_{j<k}|\ee^{\ii\theta_j}-\ee^{\ii\theta_k}|^2,
 \qquad 0\le\theta_j<2\pi.
\end{equation}
This is a familiar determinant-product example
\cite[Section 4]{Borodin2009DPP}; its relevant calculation is short.
Put $\phi_m(\theta)=\ee^{\ii m\theta}/\sqrt{2\pi}$ for $0\le m<N$.
The Vandermonde identity writes \eqref{eq:cue-joint} as
$|\det[\phi_m(\theta_j)]|^2/N!$. On expanding both determinants,
orthonormality kills all unequal permutations after full integration,
so the density has mass one. If only $N-k$ coordinates are integrated,
orthogonality instead retains the $k\times k$ minors; expanding their
determinants gives
$\rho_k=\det[K_N(\theta_a,\theta_b)]_{a,b\le k}$, where
\[
 K_N(\theta,\phi)=\frac1{2\pi}\sum_{m=0}^{N-1}
                                      \ee^{\ii m(\theta-\phi)}.
\]
In particular $\rho_2=(N/(2\pi))^2-|K_N|^2$.
Unfolding by $x=N\theta/(2\pi)$ and summing the geometric series yields
\begin{equation}\label{eq:cue-sine}
 \rho_{2,N}^{\rm unf}(s)
    =1-\left(\frac{\sin\pi s}{N\sin(\pi s/N)}\right)^2
       \longrightarrow 1-\left(\frac{\sin\pi s}{\pi s}\right)^2
\end{equation}
for fixed $s$, with removable values interpreted by continuity.
If each coordinate of the limiting real-line pair density is smoothed
by $g_\epsilon$, its separation profile becomes
\begin{equation}\label{eq:smoothed-sine-reference}
 1-g_{\sqrt2\epsilon}*\left(\frac{\sin\pi s}{\pi s}\right)^2.
\end{equation}
This follows from the same difference-variance calculation as
Proposition~\ref{prop:separation}. These equations provide a reference
for what a model marginal calculation can produce. Neither the
Vandermonde law nor its sine profile has been derived for the localized
canonical zero configuration in this manuscript.
